% Sections only — \chapter{Related Work}\label{ch:related} is in main.tex.

\section{Clearing, Clients and Central Counterparties}\label{sec:rw-clearing}

 \cite{menkveld_vuillemey_2021} offer a systematic review of the theories and empirical studies on Central Clearing. They highlight a research questions that remain open for further work -- the workings of the client clearing market. Commodity Futures Trading Commission (CFTC) data indicates that as of January 2025, clients make up 73\% of margin posted at major CCPs, but despite this studies on client clearing remain scarce. \cite{ofr_2026} study client clearing and its implications for market structure and stability by looking at Credit Default Swaps. They find that the modal client has only a single clearing connection. In periods of stress, these dealers may also stop the provision of clearing services - cutting off their clients' access to CCPs. Based on this, the simplified ABM assigns each client to a single CM who can freeze the provision of clearing services.

\cite{galbiati_soramaki_2013} build a simplified and stylised network model of tiered CCP-clearer-client relationships. They find that tiering reduces CCP total exposure but concentrates it onto fewer, larger GCMs. These extreme exposures to the CCP become larger and less predictable, despite a reduction in the average large exposure. 
 \cite{borovkova_2013} find that while CCPs reduce systemic risk in homogeneous systems, the effect depends on the system's topology and degree of homogeneity. The model here is also heterogenous, clients and CMs vary in balance sheet sizes, anchored to realistic ranges. The assignment is relatively heterogenous, with some CMs carrying upto sixteen members and others carrying as little as three.


Contagion in the model is transmitted through a well documented mechanism. In stressful periods, leveraged institutions need to liquidate in volatile markets to meet margin calls and maintain capital ratios, leading to a procyclical spiral of further calls. \cite{thurner_2012} show that heavy tailed returns and clustered volatility arise from leverage and leverage limits and the liquidation of assets due to these constraints. \cite{aymanns_farmer_2015} further look at the dynamics of these leverage cycles, and find that countercyclical leverage rules dampen the cycle. \cite{bookstaber_2014} make use of an Agent-based model to study the effects of asset and funding based fire sales on financial vulnerability. Together, these papers establish the base of how contagion transmits through the model. CMs and clients have capital  constrains and may need to liquidate during stressful periods. 

\cite{haynes_mcphail_zhu_2019} further state that the Basel III leverage ratio serves as a likely binding constraint for derivative positions. \cite{acosta_smith_2018} further show that leverage ratio affected dealers reduce client clearing, restricting clients' access to central clearing. These elements motivate the capital constraints imposed on Bank and Non-Bank Clearing Members and their clients in the Agent-Based model, and how breaching these leads to a freeze of trading behaviour and client clearing services. 

\subsection{Margins and Procyclicality}

Through the advent of central clearing, risk does not disappear; counterparty risk is rather transformed into liquidity risk \citep{cont_2017}. Accounting losses in a bilateral system with Mark-to-Market changes turn into realised margin calls drawing on liquidity in central clearing. As risk-based margins rise with volatility, these margin calls are pro-cyclical \citep{brunnermeier_pedersen_2009} -- margins peak when funding is low, leading to a spiral which causes forced sales that further depress prices and lead to even more margin calls. \cite{glasserman_wu_2018} find that volatility spikes lead to destabilising margin calls in times of market stress. \cite{esrb_2020} report that at the start of the COVID-19 crisis, IM posted at the four largest EU/UK CCPs rose by a third of pre-crisis levels, daily VM calls at the same time quintupled on euro area investment funds' derivative exposures, and increases at this time stemmed mostly from client portfolios and to a lesser extent from house portfolios. Prompted by the large March 2020 margin calls, \cite{bcbs_cpmi_iosco_2025} set out to review CCP margining practices and created a policy proposal calling for greater transparency in margining practices by CCPs onto CMs and further onto clients. The stressful period used here is also calibrated on the COVID-19 onset data. 

 \cite{paddrik_rajan_young_2020} use CDS trade level data to model how a credit shock triggers margin calls through a network of members and clients. They find that large client firms in the periphery with imbalanced portfolios contribute most to contagion. Through this model, we also see that a few large client defaults overwhelm the CMs more than many small defaults. An alternative to margins that react to recent market volatility are "through the cycle" margin models that ask for higher margins in quiet times \citep{glasserman_wu_2018}. These additional requirements, also called anti-procyclicality tools, minimise destabilising increases in IM requirements during market stress \citep{blanck_2025}. We therefore look at how these constant through the cycle margins compare with reactive VaR based margining as a demonstrative test for our model.

\section{Agent-Based Models}

Analysis of CCP risk and practises can only go so far with historical data. Agent-Based modelling is a novel methodology to study complex, adaptive and emergent behaviours arising from the interactions of heterogeneous agents \citep{axtell_farmer_2025}, and through this it offers an excellent approach for the analysis of central clearing. The analysis of these behaviours is particularly useful when looking at stressed market environments, and ABMs have found increasing popularity in recent years in simulating financial markets. \cite{deloitte_ccp} develop an agent-based CCP risk model to analyse the effects of proposed recommendations to CCP risk practices laid out in a white paper by major global CMs \citep{ccp_path_forward_2020}, such as the adequacy of Cover-2 for CCP default funds. They incorporate a simulated market environment with an exchange and CMs, complete with a limit order book where CMs submit orders and exhibit behaviours found in fundamental, momentum and noise traders. The CCP then collects margins based on these trades, which each CM needs to answer, alongside a default management framework to simulate realistic CCP behaviour. The model in this paper is in many ways an extension of this model, adopting a similar architecture and extending it with a client tier.\footnote{The publicly available descriptions of this model are the
practitioner white paper \citep{deloitte_ccp} and Simudyne's online model
documentation; the detailed model documentation on which
Chapter~\ref{ch:model} draws was provided to the author by Simudyne and is cited as unpublished technical documentation \citep{simudyne_odd}.}

\subsection{Financial Market Simulations}

The first stage in the model is the design of a simulated market environment. The market simulations attempt to reproduce the stylised facts of returns we see in the market -- heavy tails, near zero autocorrelations and volatility clustering. One approach behind this is to look at price dynamics through linear price impact caused by the aggregate and excess demands of different trader agents. An early paper on this is \cite{chiarella_1992}, which populates the market with fundamentalists trading on a fundamental value and chartists trading on momentum and trend. An extension of this model is proposed by \cite{majewski_2018}, who add noise traders and allow for a long term drift in fundamental asset value and create an Agent-Based Model (Extended-Chiarella Model). In their model, the parameters carried by each trader are jointly estimated alongside the Fundamental Value using a Kalman Filter and an Expectation Maximisation algorithm. This offers inspiration for this paper, which calibrates the parameters of a synthetic Fundamental Value signal alongside agent parameters. 

Another approach replaces the aggregate price impact mechanism with an explicit limit order book (LOB) such that prices emerge from the matching of submitted orders. \cite{chiarella_iori_perello_2009} place fundamentalists, chartists and noise traders in such a system. \cite{farmer_2005} show that noise (Zero-Intelligence) traders submitting orders - limit orders, market orders, and cancellations - at random in a continuous double auction are adequate in explaining 96\% of variance in the bid-ask spread and 76\% of the variance of the price diffusion rate. \cite{cont_stoikov_talreja_2008} propose a stylised model for the dynamics of a limit order book, where order flows are governed by independent Poisson processes, separated into independent rates for market orders, limit orders and order cancellations. \cite{gao_2022_hfabm} construct a High Frequency ABM at millisecond levels to simulate the 2010 Flash Crash, and use a surrogate modelling technique to calibrate parameters.

The financial simulation in this paper is based upon these models, with a discrete time LOB operating at a 1-min frequency, populated by Fundamental Traders, Momentum Traders and Noise Traders. Fundamental Traders trade according to an exogenous fundamental value signal, Momentum Traders trade on a momentum signal and Noise Traders place orders at independent rates for market orders, limit orders and order cancellation. The model is calibrated on real market data and includes a calm and a stressed period to investigate resilience and margin procyclicality under stress.



\subsection{Model Calibration}

The credibility of an agent-based model is dependent on how well it is calibrated, i.e. how well it captures realistic behaviours. For the market simulation, the calibration approach used is a minimisation of distance from stylised facts of returns we see on the market. \cite{cont_2001} present a set of stylised facts for returns such as heavy tails, an absence of correlation, volatility clustering and a slow decay of autocorrelation in absolute returns. \cite{cont_2007} find that for agent-based models, volatility clustering in particular remains difficult to reproduce, and the ones that manage to do so include a switching of the market between high and low activity periods. 

To measure distance from these stylised facts, a loss function must first be defined that quantifies this robustly. \cite{franke_westerhoff_2012} calibrate a financial market simulation agent-based model by minimising the distance between simulated and empirical moments, by weighing each moment by the inverse of its sampling variability so that more weight is put onto less noisy moments. They construct a sample of these moments using a block-bootstrap approach on empirical data. They then validate these moments using moment specific p-values and a moment coverage ratio. Calibration of these moments can be computationally very expensive, as they require a full simulation of the model. \cite{lamperti_2018} address this by pioneering a surrogate modelling approach for the calibration of agent-based models. They train a machine learning surrogate - in this case an XGBoost model - to approximate parameter-to-loss mapping and use active learning to concentrate around the optimum. 

\cite{gao_2022_xgb} further extend this to an XGB-Chiarella approach to estimate the parameters of a Chiarella model. An XGBoost model is trained on a random sample of parameters to predict a loss function. \cite{gao_2022_hfabm} use this methodology to identify a promising region of parameters for their High Frequency Agent-Based market simulation and then perform a grid search around the region to obtain optimal parameters. The thesis makes use of this machine learning surrogate modelling approach to calibrate parameters and steer model design. We also apply this approach to jointly calibrate the parameters of a synthetic FV signal alongside the other parameters, a novel application that combines this calibration methodology with a precedence for joint calibration through \cite{majewski_2018}.